% Data Analyst --- pengalaman (ATS, ID)
\role{IT Supervisor \& Data Analyst}{PT.\ Yih Quan Footwear Indonesia}{Feb 2025 -- Sep 2025}{Indonesia}\begin{itemize}
  \item Mengawasi dan memastikan pelaksanaan Quality Control (QC) dalam proses operasional manufaktur footwear.
  \item Melakukan monitoring dan evaluasi hasil QC untuk mengidentifikasi ketidaksesuaian, tren kualitas, dan permasalahan operasional.
  \item Mengolah hasil pemeriksaan dan menyusun laporan operasional/QC secara berkala.
  \item Menyajikan hasil analisis dan laporan kepada stakeholder terkait sebagai bahan evaluasi dan pengambilan keputusan.
  \item Berkoordinasi dengan tim produksi dan pihak terkait dalam menindaklanjuti temuan QC dan permasalahan operasional.
  \item Memastikan informasi dan hasil monitoring tersampaikan secara akurat kepada stakeholder.
  \item Mendukung proses evaluasi dan perbaikan kualitas berdasarkan hasil monitoring dan data operasional.
  \item Menggunakan Microsoft Excel untuk pengolahan data QC, rekapitulasi hasil pemeriksaan, monitoring, dan penyusunan laporan.
  \item Menganalisis tren hasil QC dan menyajikannya dalam bentuk laporan/grafik untuk memudahkan stakeholder dalam memonitor performa kualitas.
  \item Menyediakan live reporting hasil monitoring dan QC kepada stakeholder melalui Telegram.
\end{itemize}
\vspace{4pt}
\role{Head of Software Engineering}{PT.\ Lingkar Kreasi Teknologi}{Nov 2022 -- Jan 2025}{Tangerang, Indonesia}
\begin{itemize}
  \item Mendefinisikan metrik delivery/kualitas yang meningkatkan produktivitas tim 20\%.
  \item Menggunakan analitik traffic/biaya untuk menurunkan belanja server 30\% di tengah pertumbuhan traffic 200\%.
  \item Melacak KPI defect; quality gate mengurangi bug kritis 40\%.
  \item Melaporkan hasil terukur kepada pemangku kepentingan (kepuasan klien 95\%).
\end{itemize}
\vspace{4pt}
\role{VP DevOps Engineer}{PT.\ Lingkar Kreasi}{Nov 2019 -- Aug 2022}{Bandung, Indonesia}
\begin{itemize}
  \item Membangun metrik deployment/reliability ke dalam CI/CD (-50\% kegagalan; hari -> menit).
  \item Menganalisis utilisasi infrastruktur untuk menurunkan biaya 20\% dengan Kubernetes/Docker.
  \item Memantau ketersediaan terhadap target uptime 99.99\%.
\end{itemize}
\vspace{4pt}
\role{Backend Developer}{PT.\ Rakhasa Artha Wisesa}{Feb 2021 -- Nov 2022}{Jakarta Barat, Indonesia}
\begin{itemize}
  \item Mengoptimasi SQL menggunakan metrik query (-40\% latensi).
  \item Mengukur performa API (-50\% waktu respons).
  \item Melacak tingkat keberhasilan pembayaran (+20\% transaksi berhasil).
\end{itemize}
\vspace{4pt}
\role{Project Lead Developer}{PT.\ Lingkar Kreasi}{Nov 2018 -- Nov 2019}{Bandung, Indonesia}
\begin{itemize}
  \item Melacak KPI delivery (95\% tepat waktu, kepuasan 90\%+).
  \item Menggunakan metrik anggaran/scope untuk penghematan 15\% dan penurunan scope creep 25\%.
\end{itemize}
\vspace{4pt}
\role{Web Developer}{PT.\ Danadipa Central Niaga}{Nov 2017 -- Nov 2018}{Yogyakarta, Indonesia}
\begin{itemize}
  \item Menganalisis data performa situs untuk memperbaiki waktu muat dan metrik SEO.
\end{itemize}